% fig-i-32.tex — I.32: triangle angle sum (parallel through apex).
\begin{figure}[H]
\centering
\begin{tikzpicture}[scale=1.2, line cap=round]
  \coordinate (A) at (-1.5, 0);
  \coordinate (B) at (2.0, 0);
  \coordinate (C) at (0.5, 2.4);
  \coordinate (D) at ($(C)!-1.4!(A)$); % line through C parallel to AB, on left
  \coordinate (E) at ($(C)!-1.4!(B)$); % line through C parallel to AB, on right
  % Triangle.
  \draw[very thick] (A) -- (B) -- (C) -- cycle;
  % Parallel through C, drawn long.
  \draw[thin] ($(C)!-0.7!(B)$) -- ($(C)!1.5!(B)$);
  % Side AB extended beyond B to F to expose exterior angle.
  \coordinate (F) at ($(A)!1.4!(B)$);
  \draw[thin] (B) -- (F);
  % Labels.
  \node[below left]  at (A) {$A$};
  \node[below right] at (B) {$B$};
  \node[above]       at (C) {$C$};
  \node[below right] at (F) {$F$};
  \node[above left]  at ($(C)!-0.7!(B)$) {$D$};
  \node[above right] at ($(C)!1.5!(B)$) {$E$};
\end{tikzpicture}
\caption{Proposition I.32. Drawing $DE$ through $C$ parallel to $AB$
makes $\angle DCA = \angle CAB$ (alternate, I.29) and $\angle ECB =
\angle ABC$ (alternate, I.29); the straight angle at $C$ then sums
the three interior angles of $\triangle ABC$ to two right angles.}
\label{fig:I.32}
\end{figure}
